\chapter{Software Design}
\label{Chapter 4}

The previous chapter described the platform from the angle of the features that got built. This chapter steps back to the design level: what the design was being tested against, how the system decomposes into components, the data model that ties them together, the principal sequence flow, and the security model.

\section{Design Goals}

When the platform was being designed, a fairly short list of goals -- in priority order -- was what every architectural choice got tested against. Tenant isolation sits at the top of that list. The reason: a cross-tenant data leak in a regulated space is the kind of mistake you don't get to walk back. Configurable collection policy ranks higher than cost-aware channels for an unromantic reason: every tenant whose policy ends up hard-coded into the orchestrator eventually needs a code change just to be onboarded, and that path does not scale past about three tenants. The rest of the goals are operational in nature -- auditability, observability, graceful degradation. Table~\ref{tab:design-goals} lists the full set.

\begin{table}[H]
\centering
\caption{Prioritised design goals.}
\label{tab:design-goals}
\begin{tabularx}{\textwidth}{l X}
\toprule
\textbf{Goal} & \textbf{Design implication} \\
\midrule
Tenant isolation by default & Per-tenant identity realm, per-tenant search partition, tenant-scoped database queries, thread-local tenant context propagated across asynchronous boundaries. \\
Configurable collection policy & Workflow stored as data (nodes plus boolean transitions), not as code. Tenant edits do not require a deployment. \\
Multi-channel orchestration & Channels are workflow action types. Adding a new channel reduces to adding a new node type, not rewriting the orchestrator. \\
Auditability & Every state transition writes a row to the case timeline; every outbound action carries an idempotency key. \\
Graceful degradation & Database-backed job queue, per-job retry policy with exponential backoff, dead-letter handling for permanent failures. \\
Cost-aware channels & SIM-based WhatsApp and SMS rather than vendor-priced APIs; AI calls only when the workflow asks for them. \\
\bottomrule
\end{tabularx}
\end{table}

\section{Component Decomposition}

The platform splits into seven major components, each owning a clear slice of the system. Inter-component traffic is exclusively HTTP, authenticated with either an API key or a user token. Inside a component, the layering is conventional and is kept honest through code review rather than enforced by package boundaries. Table~\ref{tab:components} lists the components.

\begin{table}[H]
\centering
\caption{Major components of the platform.}
\label{tab:components}
\begin{tabularx}{\textwidth}{l X}
\toprule
\textbf{Component} & \textbf{Responsibility} \\
\midrule
Orchestrator backend & Cases, workflows, campaigns, scheduled jobs, inbound callbacks -- all of it. \\
Agent dashboard     & The browser UI that human agents and campaign managers actually use. \\
Automation service  & Drives the WhatsApp sender pool. Bridges the orchestrator to the multi-device library. \\
Desktop client      & Operator-side app. Locally-paired WhatsApp workers. End-of-day scraping. \\
Android sender app  & Pushes SMS out from a real consumer SIM. Pulls inbox replies every night. \\
Identity layer      & SSO. Realm per tenant. Federation into the lender's own identity provider. \\
Search and analytics & Per-tenant search over cases. Event store feeding the dashboard and the data team. \\
\bottomrule
\end{tabularx}
\end{table}

\section{Data Model}

The data model is a normalised relational schema, with denormalised projections cached in a search index for queries and in an analytics event store for reporting. The most important relationship in the model is between the case and its per-cycle allocation. The case is a stable handle on the customer's relationship with the lender. The allocation is the cycle-specific projection the workflow engine actually drives. Splitting the two keeps the case timeline append-only across cycles, and lets a case carry its history forward even when its workflow position is reset at the start of a new cycle. The second important design choice is the append-only case-event table: every action the platform takes against a case appends a row, so timeline queries read from that table directly, while queries that need the current state read materialised columns on the case and the allocation. This separation keeps the live tables small and the timeline easy to audit.

\section{Principal Flow: AI Voice Call}

If you want to understand the architecture in one flow, look at an AI voice call. It is the most instructive one to trace because it crosses every major component and exercises the workflow engine on both the outbound and the inbound side. When the orchestrator's scheduler hits an AI-call node for a case, two jobs get scheduled at once: the call-initiate job (to run now) and a defensive status-poll job (to run five minutes later). The executor for the first job pulls the next available from-number, assembles the call request, and hands it across to the internal voice AI platform. The platform places the call via the telephony gateway. When the call ends it returns a terminal callback to us. Our callback processor writes the disposition onto the case, cancels the poll job, emits the relevant metrics, and re-enters the engine -- the workflow then evaluates the outgoing transitions of the current node and walks the case forward, all in the same request. Two things in that flow are worth calling out. One: the defensive poll job exists because the callback is best-effort. Without a polling fallback, a single dropped callback could leave a case parked forever. Two: we cancel that poll job softly, not hard. The poll job is allowed to run, observes that the case has already been advanced, and exits. Under any eventual-consistency hiccup in the database, that's a more robust pattern than a hard cancel.

\section{Security Design}

Security on the platform is laid out as independent layers, so that one layer being breached does not collapse the layers below it. Users authenticate against the identity layer using short-lived signed tokens. Every API call resolves a tenant identifier from the URL prefix, and any token whose realm does not match that tenant is rejected before anything else happens. A thread-local tenant context then carries that identity through every downstream call -- including across asynchronous boundaries. Role-based access control is enforced both on the backend and on the dashboard. Every database query is scoped to the tenant. Every search query is routed to the tenant's own partition. Every analytics query carries the tenant tag in its filter. Cross-tenant queries are simply not allowed at the data-access layer at all. Inbound callbacks are authenticated either with a signed body digest or with a monthly-rotated API key, and that secret lives in the cloud provider's secrets manager rather than in any config file. Sitting alongside all of this is a separate set of compliance controls -- call-frequency guards, do-not-call honour periods, PII access logging, recording retention -- and those apply uniformly across every tenant, regardless of how that tenant has chosen to configure its own workflow.
